% sec04_3.tex — Section 4.3: Boundary Conditions
\section{Boundary Conditions}
\label{sec:wave-bc}

The partial differential equation for a vibrating string, \eqref{eq:4.2.7} or \eqref{eq:4.2.8}, has a second-order spatial partial derivative.  We will apply one boundary condition at each end, just as we did for the one-dimensional heat equation.

The simplest boundary condition is that of a fixed end, usually fixed with zero displacement.  For example, if a string is fixed (with zero displacement) at $x = L$, then
\begin{equation}
u(L,t) = 0.
\label{eq:4.3.1}
\end{equation}
Alternatively, we might vary an end of the string in a prescribed way:
\begin{equation}
u(L,t) = f(t).
\label{eq:4.3.2}
\end{equation}
Both \eqref{eq:4.3.1} and \eqref{eq:4.3.2} are linear boundary conditions; \eqref{eq:4.3.2} is nonhomogeneous, while \eqref{eq:4.3.1} is homogeneous.

A more interesting boundary condition occurs if one end of the string is attached to a dynamical system.  Let us suppose that the left end, $x=0$, of a string is attached to a spring-mass system, as illustrated in Fig.~4.3.1.  We will insist that the motion be entirely vertical.  To accomplish this, we must envision the mass to be on a vertical track (possibly frictionless).  The track applies a horizontal force to the mass when necessary to prevent the large horizontal component of the tensile force from turning over the spring-mass system.  The string is attached to the mass so that if the position of the mass is $y(t)$, so is the position of the left end:
\begin{equation}
u(0,t) = y(t).
\label{eq:4.3.3}
\end{equation}

\begin{figure}[H]
\centering
\includegraphics[width=0.8\textwidth]{figures/ch04/fig_4_3_1.png}
\caption{Spring-mass system with a variable support attached to a stretched string.}
\label{fig:4.3.1}
\end{figure}

However, $y(t)$ is unknown and itself satisfies an ordinary differential equation determined from Newton's laws.  We assume that the spring has unstretched length $l$ and obeys Hooke's law with spring constant $k$.  To make the problem even more interesting, we let the support of the spring move in some prescribed way, $y_s(t)$.  Thus, the length of the spring is $y(t) - y_s(t)$ and the stretching of the spring is $y(t) - y_s(t) - l$.  According to Newton's law (using Hooke's law with spring constant $k$),
\[
m\frac{d^2 y}{dt^2} = -k(y(t) - y_s(t) - l) + \text{other forces on mass.}
\]
The other vertical forces on the mass are a tensile force applied by the string $T(0,t)\sin\theta(0,t)$ and a force $g(t)$ representing any other external forces on the mass.  We must be restricted to small angles, such that the tension is nearly constant, $T_0$.  In that case, the vertical component is approximately $T_0\partial u/\partial x$:
\[
T(0,t)\sin\theta(0,t) \approx T(0,t)\frac{\sin\theta(0,t)}{\cos\theta(0,t)} = T(0,t)\pd{u}{x}(0,t) \approx T_0\pd{u}{x}(0,t),
\]
since for small angles $\cos\theta \approx 1$.  In this way the boundary condition at $x=0$ for a vibrating string attached to a spring-mass system [with a variable support $y_s(t)$ and an external force $g(t)$] is
\begin{equation}
\boxed{m\frac{d^2}{dt^2}u(0,t) = -k(u(0,t) - y_s(t) - l) + T_0\pd{u}{x}(0,t) + g(t).}
\label{eq:4.3.4}
\end{equation}

Let us consider some special cases in which there are no external forces on the mass, $g(t)=0$.  If, in addition, the mass is sufficiently small so that the forces on the mass are in balance, then
\begin{equation}
T_0\pd{u}{x}(0,t) = k(u(0,t) - u_E(t)),
\label{eq:4.3.5}
\end{equation}
where $u_E(t)$ is the equilibrium position of the mass, $u_E(t) = y_s(t) + l$.  This form, known as the nonhomogeneous \textbf{elastic} boundary condition, is exactly analogous to Newton's law of cooling [with an external temperature of $u_E(t)$] for the heat equation.  If the equilibrium position of the mass coincides with the equilibrium position of the string, $u_E(t) = 0$, the homogeneous version of the elastic boundary condition results:
\begin{equation}
T_0\pd{u}{x}(0,t) = ku(0,t).
\label{eq:4.3.6}
\end{equation}
$\partial u/\partial x$ is proportional to $u$.  Since for physical reasons $T_0 > 0$ and $k > 0$, the signs in \eqref{eq:4.3.6} are prescribed.  This is the same choice of signs that occurs for Newton's law of cooling.  A diagram (Fig.~4.3.2) illustrates both the correct and incorrect choices of signs.  This figure shows that (assuming $u=0$ is an equilibrium position for both string and mass) if $u > 0$ at $x = 0$, then $\partial u/\partial x > 0$ in order to get a balance of vertical forces on the \emph{massless} spring-mass system.  A similar argument shows that there is an important sign change if the elastic boundary condition occurs at $x = L$:
\begin{equation}
T_0\pd{u}{x}(L,t) = -k(u(L,t) - u_E(t)),
\label{eq:4.3.7}
\end{equation}

\begin{figure}[H]
\centering
\includegraphics[width=0.8\textwidth]{figures/ch04/fig_4_3_2.png}
\caption{Boundary conditions for massless spring-mass system.}
\label{fig:4.3.2}
\end{figure}

the same sign change we obtained for Newton's law of cooling.

For a vibrating string, another boundary condition that can be discussed is the \textbf{free end}.  It is not literally free.  Instead, the end is attached to a frictionless vertically moving track as before and is free to move up and down.  There is no spring-mass system, nor external forces.  However, we can obtain this boundary condition by taking the limit as $k \to 0$ of either \eqref{eq:4.3.6} or \eqref{eq:4.3.7}:
\begin{equation}
T_0\pd{u}{x}(L,t) = 0.
\label{eq:4.3.8}
\end{equation}
This says that the vertical component of the tensile force must vanish at the end since there are no other vertical forces at the end.  If the vertical component did not vanish, the end would have an infinite vertical acceleration.  Boundary condition \eqref{eq:4.3.8} is exactly analogous to the insulated boundary condition for the one-dimensional heat equation.

\begin{exercises}{4.3}

\item If $m = 0$, which of the diagrams for the right end shown in Fig.~4.3.3 is possibly correct?  \emph{Briefly} explain.  Assume that the mass can move only vertically.

\begin{figure}[H]
\centering
\includegraphics[width=0.8\textwidth]{figures/ch04/fig_4_3_3.png}
\caption{}
\label{fig:4.3.3}
\end{figure}

\item Consider two vibrating strings connected at $x = L$ to a spring-mass system on a vertical frictionless track as in Fig.~4.3.4.  Assume that the spring is unstretched when the string is horizontal (the spring has a fixed support).  Also suppose that there is an external force $f(t)$ acting on the mass $m$.

\begin{enumerate}
\item[(a)] What ``jump'' conditions apply at $x = L$ relating the string on the left to the string on the right?

\item[(b)] In what situations is this mathematically analogous to perfect thermal contact?
\end{enumerate}

\begin{figure}[H]
\centering
\includegraphics[width=0.8\textwidth]{figures/ch04/fig_4_3_4.png}
\caption{}
\label{fig:4.3.4}
\end{figure}

\end{exercises}
